% ID: SEL-D0018-BD-E07
% ESTADO: validada. S019. Revisión matemática, didáctica y visual.
\textbf{a) Intervalo al 98\,\%.}\par\nopagebreak[4]
La proporción observada es $\hat p=2/3$. Hay $400$ éxitos y $200$ fracasos, de modo que la aproximación normal es adecuada bajo el modelo muestral. Con el cuantil bilateral $z_{0.99}\approx2.326347874$,
\[
\hat p\pm z_{0.99}\sqrt{\frac{\hat p(1-\hat p)}{600}}
\approx\boxed{(0.621896,\ 0.711437)}.
\]
Como $13/20=0.65$ está dentro del intervalo, ese valor es compatible con los datos al nivel de confianza solicitado.

\textbf{b) Margen.}\par\nopagebreak[4]
La semiamplitud es $z_{0.99}\sqrt{(2/9)/600}\approx\boxed{0.044771}$, aproximadamente $4.48$ puntos porcentuales.

\textbf{c) Tamaño mínimo.}\par\nopagebreak[4]
Al fijar la confianza y la proporción de planificación,
\[
z_{0.99}\sqrt{\frac{2/9}{n}}<0.02
\quad\Longleftrightarrow\quad
n>\frac{z_{0.99}^2(2/9)}{0.02^2}\approx3006.608017252.
\]
Con el cuantil preciso, $3006$ incumple la cota y $\boxed{3007}$ la satisface. Si se adopta $2.33$ de una tabla de dos decimales, la cota aproximada sube a $3016.0556$ y resulta $3017$. Conviene conservar el cuantil sin redondear hasta calcular el mínimo.
